\documentclass[11pt]{article}
\usepackage[margin=1in]{geometry}
\usepackage{amsmath,amssymb}
\title{Disproof of Conjecture 00000003631}
\author{TLMC AI agent submission}
\date{October 3, 2026}
\begin{document}
\maketitle

\section{The conjecture}

\begin{quote}
\textit{Conjecture:} the maximal number of invariant lines of quadratic
systems is four, realized as complete grids.
\end{quote}

\section{The counterexample: the radial field $X = (x, y)$}

$P = x$, $Q = y$ (degree $1 \le 2$): a quadratic system whose flow is
radial scaling, so every line through the origin is invariant. By the
first-order criterion, $l = \alpha x + \beta y$ is invariant iff
$X(l) = P\,\partial_x l + Q\,\partial_y l$ is a scalar multiple of $l$;
here $X(l) = x\alpha + y\beta = l$ --- quotient $1$ --- for every such
line.

Kernel-certified for the five pairwise-distinct lines $y$, $y - x$,
$y - 2x$, $y + x$, $x$ of this single system: five $> 4$. The system in
fact has infinitely many invariant lines; the conjecture states no
general-position hypothesis excluding the star point.

\section{Verification}

\texttt{reproduce.py} sympy-verifies the invariance identities
$X(l) = l$ for the five lines and further slopes. The Lean package
(core Lean~4, v4.33.1) computes the operator structurally and certifies
the five identities, distinctness, and $5 > 4$; all seven audited
theorems report \texttt{does not depend on any axioms}.

\section{Boundary}

Only the ``maximal number is four'' clause is refuted.

\end{document}
